\section{Discussion}\label{sec:discussion}

The results split cleanly in two, and the split is the finding.

Everything that measures the \emph{supply} of liquidity moves on expiry Thursdays. Spreads
widen, displayed notional falls, a given order moves the price further, per-trade impact
rises, realised variance per minute rises, and orders are cancelled within a second of entry
more often. Trading in the settlement window costs more on the day the settlement is struck.

Everything that measures \emph{direction} does not move. The variance ratio does not rise;
it falls slightly. Signed order flow is no more persistent. The share of price impact that
survives a minute is unchanged. Submitted flow is no more one-sided. Volume does not
concentrate into the final minute, and less of it transacts at the running benchmark.

The natural reading is that the window is crowded rather than pushed. More participants are
active, for the ordinary reason that positions have to be closed or rolled, and the book
prices that congestion. Nothing in the order flow looks like one side working steadily
against the other.

\subsection{The awkward result}

The settlement VWAP does end up further from the closing price on expiry sessions, and that
is what a price pushed and then released looks like. It is also what a noisier price series
produces around any fixed reference point, and realised variance per minute is higher on
these sessions. The two explanations are not separable from the reversal alone.

They are separable using the direction measures, and those favour the second. A push large
enough to move a thirty-minute average by several basis points would have to leave some trace
in return autocorrelation, in flow persistence, or in the permanent share of impact. None of
the three moves. Taken together the evidence points to a wider, thinner, noisier market
rather than a directed one---though this is an inference from the absence of a signature, and
absence of a signature in a sample of this size is weaker evidence than its presence would
have been.

\subsection{Concealment}

The concealment results are the ones that most needed order-level data, and they are also the
ones most easily over-read. The level of concealment does not change. What changes is that
concealed parent orders get larger relative to displayed ones while being refilled less often
per share traded. That is a description of larger institutional orders in the window, which
is exactly what expiry-related hedging would look like. It is not evidence of stealth, and
reading it as such would be reading a routine execution choice as an intention.

\subsection{Where the cost falls}

The sweep measure gives the liquidity contrast a number: moving the price by trading a fixed
notional costs several times more in the illiquid group than in the liquid one. That is the
constraint that would bind anyone trying to influence a settlement, and it binds least
exactly where the incentive is cheapest to act on. The difference-in-differences in
Table~\ref{tab:contrasts} tests whether the expiry effects are in fact larger there. With the
universe used here the contrast lacks the securities to resolve it, which is a limitation of
the sample rather than a finding.

\subsection{What would settle the question}

Three things, in order of how much they would add.

Open interest per security, which is the variable the settlement story actually predicts on.
If the expiry effects scale with the size of the expiring position, that is a far more
specific prediction than anything tested here, and failure to scale would be close to
decisive against.

A wider universe. The difference-in-differences are the right tests and they are
underpowered; a few dozen securities per group would resolve them.

The futures tape. Without it the basis cannot be computed, the lead-lag between futures and
cash cannot be tested, and the hypotheses listed in Section~\ref{sec:limitations} stay open.
